\documentclass[11pt,a4paper]{article}

\usepackage[margin=1.05in]{geometry}
\usepackage{amsmath,amssymb,amsthm}
\usepackage{mathtools}
\usepackage{stmaryrd}
\usepackage{booktabs}
\usepackage{array}
\usepackage{listings}
\usepackage{xcolor}
\usepackage{graphicx}
\usepackage[colorlinks=true,linkcolor=black,citecolor=black,urlcolor=blue]{hyperref}
\usepackage{enumitem}
\usepackage[T1]{fontenc}
\usepackage[protrusion=true,expansion=false]{microtype}
\setlength{\emergencystretch}{2.5em}

%%%%%%%%%%%%%%%%%%%%%%%%%%%%%%%%%%%%%%%%%%%%%%%%%%%%%%%%%%%%%%%%%%%%%%%%%%%%%%%
% Theorem environments
%%%%%%%%%%%%%%%%%%%%%%%%%%%%%%%%%%%%%%%%%%%%%%%%%%%%%%%%%%%%%%%%%%%%%%%%%%%%%%%
\theoremstyle{plain}
\newtheorem{theorem}{Theorem}[section]
\newtheorem{proposition}[theorem]{Proposition}
\newtheorem{lemma}[theorem]{Lemma}
\newtheorem{corollary}[theorem]{Corollary}
\newtheorem{conjecture}[theorem]{Conjecture}

\theoremstyle{definition}
\newtheorem{definition}[theorem]{Definition}
\newtheorem{example}[theorem]{Example}
\newtheorem{principle}[theorem]{Principle}
\newtheorem{obligation}[theorem]{Obligation}
\newtheorem{observation}[theorem]{Observation}
\newtheorem{prediction}[theorem]{Prediction}

\theoremstyle{remark}
\newtheorem{remark}[theorem]{Remark}

%%%%%%%%%%%%%%%%%%%%%%%%%%%%%%%%%%%%%%%%%%%%%%%%%%%%%%%%%%%%%%%%%%%%%%%%%%%%%%%
% Macros.  NOTATION RULES OBSERVED THROUGHOUT:
%   * the rho calculus is never written with the Greek letter rho
%   * quoting is @P, dropping is *x; corner quotes are pre-2016 notation
%%%%%%%%%%%%%%%%%%%%%%%%%%%%%%%%%%%%%%%%%%%%%%%%%%%%%%%%%%%%%%%%%%%%%%%%%%%%%%%
\newcommand{\rhoc}{\textsf{rho}}
\newcommand{\sem}[1]{\llbracket #1 \rrbracket}
\newcommand{\V}{\mathbb{V}}
\newcommand{\Rnn}{\mathbb{R}_{\geq 0}}
\newcommand{\Cx}{\mathbb{C}}
\newcommand{\Nat}{\mathbb{N}}
\newcommand{\Proc}{\mathcal{P}}
\newcommand{\Names}{\mathcal{N}}
\newcommand{\supp}{\mathrm{supp}}
\newcommand{\scong}{\equiv}
\newcommand{\tensorv}{\otimes}
\newcommand{\sumv}{\oplus}
\newcommand{\Cand}{\mathrm{Cand}}
\newcommand{\crisp}[1]{\lceil #1 \rceil}
\newcommand{\dok}{\textsf{do}}
\newcommand{\ink}{\textsf{in}}
\newcommand{\fork}{\textsf{for}}
\newcommand{\where}{\textsf{where}}
\newcommand{\wmap}{\mathsf{w}}
\newcommand{\prop}{\mathsf{a}}
\newcommand{\Rej}{\mathcal{R}_{\mathrm{rej}}}
\newcommand{\Ded}{\mathcal{R}_{\mathrm{ded}}}

\lstdefinestyle{rholang}{
  basicstyle=\ttfamily\small,
  keywordstyle=\bfseries,
  morekeywords={for,new,in,where,contract,match,if,else,do},
  columns=fullflexible,
  keepspaces=true,
  showstringspaces=false,
  breaklines=true,
  frame=single,
  framesep=5pt,
  rulecolor=\color{black!35},
  xleftmargin=6pt,
  aboveskip=9pt,
  belowskip=9pt
}
\lstset{style=rholang}

%%%%%%%%%%%%%%%%%%%%%%%%%%%%%%%%%%%%%%%%%%%%%%%%%%%%%%%%%%%%%%%%%%%%%%%%%%%%%%%

\title{\bfseries Reprogramming Across the Organ Boundary\\[4pt]
\large A graded \rhoc{} model of OSK partial reprogramming in hindbrain motor
neurons: dynamic weight maps for epigenetic state, the induction schedule as a
scoped $\sigma$-intervention, and a proofreading bound on achievable
selectivity}

\author{L.\,G.\ Meredith\\ \small F1R3FLY.io}

\date{F1R3FLY.io working note --- November 2026\\[2pt]
\normalsize Draft 1}

\begin{document}
\maketitle

\begin{center}
\fbox{\begin{minipage}{0.88\textwidth}\small
\textbf{Scope.} This note is a preclinical modelling analysis. It specifies a
computational model, the fragment of the calculus that hosts it, the
measurements the model requires, and the claims the model is and is not
competent to support. It is not a protocol, a dose, a delivery specification,
or a clinical recommendation, and nothing in it should be read as one.
\end{minipage}}
\end{center}

\begin{abstract}
Lu et al.\ showed that ectopic expression of \emph{Oct4}, \emph{Sox2} and
\emph{Klf4} (OSK) in murine retinal ganglion cells restores youthful DNA
methylation, promotes axon regeneration after optic nerve crush, and reverses
vision loss, with the effect requiring the demethylases TET1 and TET2
\cite{LuEtAl}. This note asks what a model of the analogous intervention in
\emph{motor} neurons must contain, under the constraint that the target cells
be hindbrain-derived, and it answers in the calculus of the companion notes:
the rho calculus with graded where-clauses over $\Rnn$
\cite{GradedV4,GradedPLN}, weighted GSLTs \cite{WeightedGSLT}, the
intervention membrane of \cite{Interventions}, and the proofreading ladder of
\cite{Proofreading}.

Four findings are reported. First, the published calibration does not
transfer, and the formalism says why in a form that is checkable rather than
rhetorical: a weight map is keyed by spatial--behavioural formulae refining a
rule's redexes, so a change of \emph{chromatin configuration} changes the key
set and not merely the values, and the partition discipline of
\cite{WeightedGSLT} renders an unclassified hindbrain redex a well-formedness
violation rather than a poor fit. Retinal ganglion cells are forebrain-lineage;
the bulbar motor pool is hindbrain-lineage; and Jokhai, Dundes et al.\ report
that the anterior and posterior neural ectoderm arise as parallel progenitors
with divergent chromatin landscapes \cite{TwoOrgans}. A second transfer gap,
from acute axotomy of a healthy neuron to chronic degeneration of a sick one,
is larger and less tractable, and belongs in the term rather than in a
coefficient.

Second, dynamic weight maps are necessary and not merely convenient. A locus
with $n$ CpG sites has $2^n$ methylation configurations; a static-rate
stochastic $\pi$ model must spend species on them, and SPiM is exactly the
static-map, channel-identity-key restriction of the weighted construction
\cite{WeightedGSLT}. Keying the map on a formula over the chromatin term
prices the same state dependence at one species.

Third, rejuvenation and dedifferentiation are not two processes with two
rates. They are candidates contending for the same OSK pool, hence members of
one contention set, and the therapeutic window is therefore a property of the
normalisation rather than of either rate --- a first-passage property of the
jump chain. Because the clause governs the jump chain and not the holding
times, an induction schedule must be \emph{reified as data} to come within
reach of a $\dok$ at all; a schedule left as an exogenous time-varying rate is
outside the construct. Vector tropism is namespace confinement, and the
schedule is a conditional $\sigma$-intervention in the dose reading of
\cite{Interventions}.

Fourth, the proofreading ladder yields a limitative result with an immediate
experimental consequence. Pulsed induction is proofreading in time, its
advantage over sustained induction scales as a power of the single-stage
discrimination factor $f$ at the contested loci, and the exponent is
independent of the number of pulses. Since energy buys confidence and never
scope \cite{Proofreading}, no schedule manufactures selectivity that the OSK
affinity landscape does not already supply. Measuring $f$ on
hPSC-derived hindbrain motor neurons is therefore a go/no-go test that
precedes animal work and costs a binding assay rather than a cohort.

Four places where the apparatus is not ready are recorded rather than
smoothed over: the pluripotency network's variable-level cycles, the
permanent-shaped normalisation of a well-mixed cell's contention sets,
structural identifiability, and the discipline required to distinguish a
biologically real starvation from an over-restricted $\dok$. We close with
what the model can and cannot license. Its job is to say which experiments to
run and which protocol variants to discard --- not which protocol to use.
\end{abstract}

\tableofcontents

%%%%%%%%%%%%%%%%%%%%%%%%%%%%%%%%%%%%%%%%%%%%%%%%%%%%%%%%%%%%%%%%%%%%%%%%%%%%%%%
\section{Introduction}\label{sec:intro}

\subsection{The question}

The optic-nerve result of \cite{LuEtAl} is the cleanest demonstration that
partial reprogramming can restore regenerative competence to a central nervous
system neuron without erasing its identity. Three factors suffice --- Oct4,
Sox2, Klf4, with c-Myc withheld on oncogenic grounds --- delivered by AAV under
a tetracycline-controlled switch, and the restoration of youthful methylation
depends on TET1 and TET2.

The question this note addresses is what must be true of a \emph{model} of the
analogous intervention in motor neurons, under one anatomical constraint: that
the cells be hindbrain-derived. The constraint is motivated by
\cite{TwoOrgans}, which reports that the brain develops from two parallel
progenitor populations rather than one, and that the posterior population ---
the hindbrain's --- had resisted derivation in vitro precisely because it is
not reachable from the anterior one.

The note does not propose an intervention. It specifies a model, locates it in
a fragment of the calculus where its semantics is exact and its cost is
bounded, derives from it the measurements it needs, and states what it is
competent to conclude. The last of these is the load-bearing contribution: a
model of a therapeutic cascade is worth having only in proportion to the
sharpness with which its own limits are stated.

\subsection{Why this calculus}

Three properties of the graded where-clause programme make it the right host,
and one of them is decisive.

The decisive one is that the central biological question is a question about
\emph{branches}. Partial reprogramming succeeds when a cell recovers
regenerative competence and fails when it loses terminal identity, and these
are competing outcomes of the same driving process. A model whose dynamics are
confluent --- a system of rate equations, a mean-field trajectory --- can
\emph{state} that the second outcome has some probability, but it cannot
\emph{inhabit} the branching that makes the therapeutic window a window. This
is the observation of \cite{Inhabiting} transposed: on a confluent witness
language the linear-time/branching-time spectrum collapses, and the structure
that a therapeutic window \emph{is} lives in the part that collapsed.

The second is that the state dependence of the kinetics is a dependence on
\emph{shape}. Methylation state is a configuration of a locus, and the
propensity of the next rewrite depends on it. A weight map keyed by
spatial--behavioural formulae over the term expresses that directly
(\S\ref{sec:layers}); a static rate attached to a channel identity cannot, and
pays for the difference in species.

The third is that the intervention has a home. The construct
$\dok(C)\ \ink\ P$ of \cite{Interventions} is a policy conditional on a race
arising, scoped at the cut, which is the right shape for a therapeutic
intervention: one does not know in advance which cell will be in which state
when the vector arrives, so one installs a policy rather than performing a
surgery on a named node.

\subsection{Contributions}

\begin{enumerate}[leftmargin=1.4em,itemsep=3pt]
\item \textbf{A transfer obstruction with a checkable form}
(\S\ref{sec:transfer}). The anterior calibration does not carry to the
posterior organ by reparameterisation, because the weight map's keys and not
only its values change; the partition discipline makes the failure a
well-formedness violation. A second obstruction --- acute injury versus
chronic degeneration --- is identified as the larger of the two.
\item \textbf{Dynamic weight maps as a necessity} (\S\ref{sec:layers}). The
$2^n$ argument, and the precise sense in which the model lies strictly beyond
SPiM.
\item \textbf{The window as a contention-set property} (\S\ref{sec:race}).
Rejuvenation and dedifferentiation contend for one OSK pool; the window is a
first-passage property of the normalised jump chain; and the schedule must be
reified as data to be intervenable.
\item \textbf{A proofreading bound and a cheap test}
(\S\ref{sec:fidelity}). Pulsing is proofreading in time; its benefit is a
power of the single-stage discrimination factor; and the scope limitation of
\cite{Proofreading} converts into a go/no-go measurement that precedes animal
work.
\item \textbf{Counterfactual cohort design} (\S\ref{sec:cohort}). Replay under
$\dok$ with the exogenous record held gives matched-pair contrasts per cell,
and the variance reduction is the practical payoff of rung three.
\item \textbf{Four gaps, and what they cost} (\S\ref{sec:gaps}), and a
statement of licensed and unlicensed conclusions (\S\ref{sec:license}).
\end{enumerate}

\subsection{What this note does not do}

It does not fit parameters; most of the relevant rate constants are not known
and \S\ref{sec:gaps} argues that fitting before an identifiability analysis
would produce a curve rather than a model. It does not treat vector
biodistribution, capsid immunogenicity, or whole-organism outcome, all of
which the model consumes as inputs. It does not discharge the cyclic case at
variable level. It does not price a $\dok$ against the ledger, which remains
Obligation~5.3 of \cite{Interventions}. And it proposes no protocol.

%%%%%%%%%%%%%%%%%%%%%%%%%%%%%%%%%%%%%%%%%%%%%%%%%%%%%%%%%%%%%%%%%%%%%%%%%%%%%%%
\section{Scope and substrate}\label{sec:scope}

\subsection{The target compartment}

The modelled substrate is the hindbrain-derived bulbar motor pool: the cranial
motor nuclei --- hypoglossal, facial, trigeminal motor, and the
vagal/ambiguus column --- arising from the \emph{Gbx2}-expressing posterior
neural ectoderm of \cite{TwoOrgans}. The choice is not presentational. It
fixes the weight maps, and \S\ref{sec:transfer} argues it blocks reuse of the
published calibration.

\subsection{Two populations, one delivery term}

The constraint has a consequence that must be built in rather than noted. The
medullary respiratory nuclei occupy the same compartment, so a model
containing only the intended target population cannot represent the dominant
risk at all: there is no term in it for the thing one is worried about.

\begin{principle}[Off-target population is a modelled population]
\label{prin:offtarget}
The model carries at least two cell populations under one delivery term: the
intended bulbar motor pool, and a neighbouring population receiving the same
vector exposure and the same induction schedule but differing in chromatin
state and in the consequence of mis-reprogramming. A model with one population
and a scalar ``off-target rate'' is not a model of the risk; it is a parameter
named after it.
\end{principle}

The relevance of the compartment is not incidental to the disease.
\cite{TwoOrgans} notes that amyotrophic lateral sclerosis affects both organs,
so a hindbrain-restricted intervention is partial by construction, and the
model should report what fraction of the motor deficit its target compartment
can in principle address rather than leaving that fraction implicit.

\subsection{Deliberate exclusions}

Four things lie outside the model: in vivo vector biodistribution and
pharmacokinetics; the immune response to the capsid; whole-organism functional
outcome; and anything resembling a human dose. The first three are measurement
problems the model consumes as inputs. The fourth is not a modelling question
and does not become one by being asked of a model.

The modelled object is therefore one cell's cascade, instantiated across a
population of two cell types, under a shared intervention term.

%%%%%%%%%%%%%%%%%%%%%%%%%%%%%%%%%%%%%%%%%%%%%%%%%%%%%%%%%%%%%%%%%%%%%%%%%%%%%%%
\section{The transfer obstruction}\label{sec:transfer}

\subsection{Keys, not values}

The retina arises from the optic vesicle, an evagination of the diencephalon,
so retinal ganglion cells are forebrain-lineage neurons. The result of
\cite{LuEtAl} therefore sits entirely within the anterior organ of
\cite{TwoOrgans}, while the target of this programme sits in the posterior one.

In the weighted-GSLT setting this is not a change of parameters, and the
reason is structural.

\begin{observation}[A lineage change is a change of key set]\label{obs:keys}
A weight map $\wmap$ assigns weights to keys that are spatial--behavioural
formulae refining a rule's left-hand side \cite{WeightedGSLT}. The keys are
predicates on the \emph{shape} of the state. A reparameterisation changes
$\wmap(\varphi)$ at fixed $\varphi$. A change of chromatin configuration
changes which refinements are distinguishable --- which $\varphi$ there are ---
and so changes the domain of $\wmap$, not its values.
\end{observation}

The force of this is that the partition discipline makes the failure
detectable rather than silent. The keys of a rule's map must be pairwise
exclusive and jointly exhaustive over that rule's redexes; this single
requirement is what makes the weighting a total, single-valued classification
\cite{WeightedGSLT}. Hence:

\begin{proposition}[Unclassified redexes are a well-formedness violation]
\label{prop:wf}
Let $\wmap^{A}$ be a weight map whose partition was estimated on anterior
substrate, and let $P^{H}$ be a term over hindbrain substrate. If some redex
of $P^{H}$ satisfies no key of $\wmap^{A}$, then $\wmap^{A}$ is not a weight
map for $P^{H}$: the classification is not exhaustive, and no choice of values
repairs it.
\end{proposition}

So porting the anterior calibration and re-estimating rates has two possible
outcomes, and the discipline is what distinguishes them: either the port
throws, in which case the transfer obstruction is established, or it does not,
in which case the anterior partition is exhaustive for the hindbrain term and
the obstruction is confined to the values. Either answer is informative, and
the test is cheap because it is static.

\subsection{The larger gap: acute injury versus chronic disease}

The lineage gap is the one the formalism makes crisp. The injury-regime gap is
probably the one that decides whether the approach works, and it has had less
attention.

The published result is axon regeneration after acute optic nerve crush in an
otherwise healthy neuron, together with restored function in aged and
glaucomatous eyes. Amyotrophic lateral sclerosis is chronic, progressive
degeneration of a cell under sustained proteostatic and excitotoxic stress,
with the neuromuscular junction failing distally. ``Reprogramme a healthy cell
whose axon was severed'' and ``reprogramme a sick cell that is dying slowly''
are different computational problems.

\begin{principle}[The disease belongs in the term]\label{prin:disease}
A model that represents the diseased baseline as a reduced regeneration rate
will overstate the intervention's reach, because it deletes the competition.
The pathological processes are contending candidates at the same sites the
intervention acts on --- they consume the same proteostatic capacity and the
same transcriptional machinery --- so they belong in the parallel composition,
not in a coefficient.
\end{principle}

This is the point at which the branching argument of \S\ref{sec:intro} does
real work rather than decorative work. In a rate-equation model the disease
and the therapy are additive terms on one trajectory. In the term model they
are candidates in contention, and the quantity of interest is which fires.

\subsection{The empirical assumption this rests on}

The argument of \S\ref{sec:transfer}.1 rests on one claim the source does not
establish: that the divergent chromatin landscape of the two progenitor
populations \emph{persists into the mature motor neuron}. \cite{TwoOrgans} is
a result about progenitors at gastrulation. If the divergence largely washes
out by maturity, the lineage objection weakens substantially and the anterior
calibration becomes a defensible prior.

The two-organ claim is also recent and contested at its edges; the
distinctness of the two progenitor populations has been questioned in
commentary on the paper. The model should therefore carry lineage divergence as
a quantity with a magnitude and an uncertainty rather than as a binary, and
that magnitude is among the first things worth measuring, since the whole
question of what may be reused turns on it.

\begin{obligation}[Persistence of divergence]\label{ob:persist}
Determine whether the chromatin divergence reported between anterior and
posterior neural ectoderm persists in terminally differentiated motor neurons,
and at which loci. The deliverable is not a yes/no but a per-locus magnitude,
because the weight map's keys are per-locus predicates.
\end{obligation}

%%%%%%%%%%%%%%%%%%%%%%%%%%%%%%%%%%%%%%%%%%%%%%%%%%%%%%%%%%%%%%%%%%%%%%%%%%%%%%%
\section{The model, in four layers}\label{sec:layers}

The model is built in four layers, each of which is a known fragment of the
programme. The layering is not stylistic: each layer has a different cost
profile and a different fragment discipline, and conflating them is how such
models become intractable.

\subsection{Layer 0: species, reactions, propensity}

The base layer is ordinary stochastic chemistry: AAV particles, tTA and TRE-OSK
transcripts, OSK protein, TET1/TET2, the methylation state of the modelled
loci, and the pathological species of Principle~\ref{prin:disease}. Receipts
are persistent for-comprehensions, propensity is
$\prop = \wmap(\varphi)\cdot\sum_{k}\gamma(k)$ over classified, funded redex
positions, and the simulator is Gillespie \cite{Gillespie,WeightedGSLT}.

\begin{remark}[SPiM is the baseline, and it is free]\label{rem:spim}
Restricted to the $\pi$-calculus with static maps and channel-identity keys,
the weighted construction is \emph{exactly} the Stochastic Pi Machine
\cite{WeightedGSLT,PhillipsCardelli}. Layer~0 is therefore a SPiM model in the
precise sense, and that is a deliberate choice: it is the layer that
collaborators outside this programme can read, check and run in their own
tooling, and it costs nothing to expose because the theorem already identifies
it as a restriction of what is being built.
\end{remark}

\subsection{Layer 1: dynamic weight maps for epigenetic state}

Layer~0 cannot express the dependence that matters. The propensity of
demethylation at a locus depends on the current methylation configuration of
that locus and of its neighbours --- cooperativity, spreading, CpG-island
boundary effects. In a static-rate model this dependence must be spent in
species.

\begin{observation}[The $2^n$ argument]\label{obs:blowup}
A locus with $n$ CpG sites has $2^n$ methylation configurations. A model whose
rates are attached to channel identities and are static must introduce one
species (or one equivalence class of species, chosen in advance) per
configuration on which a rate depends. For a promoter region with twenty CpG
sites this is of order $10^{6}$ species. A weight map keyed on a
spatial--behavioural formula over the chromatin term --- ``at least $k$ of the
sites in this window are unmethylated'', ``this boundary element is
unoccupied'' --- expresses the same dependence with one species, because the
configuration is in the term and the key is a predicate on it.
\end{observation}

This is where ``dynamic maps take the construction strictly beyond SPiM''
\cite{WeightedGSLT} stops being a slogan about expressiveness and becomes the
reason the model is writable at all. Weight-map instances are part of the
state specification, so a rewrite may update them: methylation changes the map
that governs subsequent methylation, which is what an epigenetic model is.

\begin{obligation}[The key set is the modelling work]\label{ob:keys}
Choosing the refinement formulae --- which shape predicates the model
distinguishes --- is the substantive modelling decision, and it is where
domain knowledge enters. The partition discipline constrains it (exclusive,
exhaustive) but does not determine it. An over-fine partition is
unidentifiable (\S\ref{sec:gaps}); an over-coarse one deletes the mechanism.
\end{obligation}

\subsection{Layer 2: the intervention}

The intervention has three parts, and each lands in a different place in
\cite{Interventions}.

\paragraph{Tropism is namespace confinement.} The vector reaches the cells it
reaches. In the calculus this is discipline $\mathcal{T}_4$: every constrained
receipt lies within the deployer's seal, so the $\dok$ reaches only the code
its author paid for. The capability to intervene is the capability to deliver,
and Proposition~2.3 of \cite{Interventions} --- that no context hard-intervenes
over $\Rnn$ --- is the formal counterpart of the fact that an untransduced cell
is not reprogrammed.

\paragraph{The schedule is a conditional $\sigma$-intervention.} The induction
regime is not a hard $\dok$. It is a graded, conditional, scoped restriction:
\emph{if} a race for OSK occupancy arises in a transduced cell, resolve it
with weights set by the current state of the switch. In the taxonomy of
\cite{Interventions} this is a conditional $\sigma$-intervention with
constraint in $\mathcal{F}_1$, and its induced parent set is read off the free
names of the constraint --- which is to say, the schedule's dependencies are
written in the schedule.

\begin{lstlisting}
do( occ |= dose(*clock) ) in { cell_1 | ... | cell_m | clock }
\end{lstlisting}

\paragraph{The schedule must be reified as data.} This is the non-obvious
requirement, and it is forced.

\begin{proposition}[An exogenous schedule is outside the construct]
\label{prop:reify}
The clause governs the jump chain and not the holding times
\cite[Rem.~3.2]{GradedPLN}, so an intervention on \emph{when} a site fires is
not expressible as a clause restriction. An induction schedule is precisely an
intervention on when. Hence a schedule supplied as an exogenous time-varying
rate lies outside the reach of any $\dok$, and the do-calculus results of
\cite{Interventions} do not apply to it. Reifying the schedule as a process
emitting tokens on a clock channel returns it to the clause's reach, because
the consumption of a clock token is itself a race.
\end{proposition}

\begin{proof}
Immediate from the cited remark together with
Definition~3.2 of \cite{Interventions}: the membrane reads
$\psi_r \tensorv \crisp{\varphi}$ at receipts binding constrained variables,
and a holding time is not a candidate at such a receipt. Once the schedule is
a message on a channel, the receipt that consumes it binds a variable, and the
variable is constrainable.
\end{proof}

The practical content of Proposition~\ref{prop:reify} is that schedules become
\emph{terms}, so comparing two induction regimes is comparing two terms under
the same model, and the comparison is an interventional contrast rather than
two separate simulations with different parameter files. That is also what
makes the identifiability question of \S\ref{sec:gaps} well posed.

\subsection{Layer 3: fidelity}

The fourth layer is the discrimination between intended and unintended OSK
targets, treated as a kinetic proofreading problem. It is developed in
\S\ref{sec:fidelity}, because its content is a limitative result rather than a
modelling convenience.

\subsection{Where each layer runs}

\begin{table}[t]
\centering\footnotesize
\begin{tabular}{@{}>{\raggedright\arraybackslash}p{0.13\textwidth}>{\raggedright\arraybackslash}p{0.30\textwidth}>{\raggedright\arraybackslash}p{0.22\textwidth}>{\raggedright\arraybackslash}p{0.25\textwidth}@{}}
\toprule
\textbf{layer} & \textbf{what it carries} & \textbf{fragment} & \textbf{cost}\\
\midrule
L0 & species, reactions, propensity & static maps, channel keys (= SPiM) & Gillespie step\\
L1 & epigenetic state dependence & dynamic maps, $\mathcal{F}_1$ keys & $O(2^{d})$ per candidate, $d$ = formula depth\\
L2 & tropism, schedule, dose & $\mathcal{T}_4$, $\mathcal{F}_1$, schedule reified & one COMM plus a conjunct\\
L3 & locus discrimination & base rules, funded & tokens per discard step\\
\bottomrule
\end{tabular}
\caption{The layers, their fragments, and what each costs. Fragment and
discipline names are those of \cite{Interventions}.}\label{tab:layers}
\end{table}

Table~\ref{tab:layers} records the architectural decision that
\S\ref{sec:gaps} will defend: the $\dok$ membrane is confined to Layer~2's
designated sites, and the bulk chemistry of Layer~0 runs underneath it without
one. Putting the whole cell under a $\dok$ is the mistake that makes the model
intractable.

%%%%%%%%%%%%%%%%%%%%%%%%%%%%%%%%%%%%%%%%%%%%%%%%%%%%%%%%%%%%%%%%%%%%%%%%%%%%%%%
\section{The central race}\label{sec:race}

\subsection{Rejuvenation and dedifferentiation contend}

Let $\Rej$ be the rule class that demethylates regeneration-competence loci,
restoring the youthful pattern, and $\Ded$ the class that activates the
pluripotency network and costs the cell its terminal identity. Both are driven
by the same OSK pool.

\begin{observation}[One contention set, not two rates]\label{obs:contend}
Two candidates contend when they share a message or share a linear receipt
\cite[Def.~2.4]{Interventions}. A candidate in $\Rej$ and a candidate in
$\Ded$ that consume the same OSK molecule share a message, hence lie in one
connected component of the contention relation, hence in one contention set
--- which is the unit at which weights are normalised.
\end{observation}

The consequence is the structural claim of this note.

\begin{proposition}[The window is a normalisation property]\label{prop:window}
Under Observation~\ref{obs:contend} the probability that a cell reaches the
rejuvenation threshold before the dedifferentiation threshold is a
first-passage property of the jump chain whose transition probabilities are
the contention-set-normalised clause values, and it is not a function of the
two rule classes' absolute propensities separately. Scaling both by a common
factor leaves it unchanged; scaling either alone changes it.
\end{proposition}

\begin{proof}
Resolution normalises within a contention set
\cite[Def.~2.6]{Interventions}: the probability that a maximal matching fires
is its weight divided by the total over the set. A common positive scaling of
all weights in a set leaves every ratio fixed, hence leaves the jump chain
fixed; first passage is a property of the jump chain alone, the holding times
being irrelevant to which threshold is reached first.
\end{proof}

Proposition~\ref{prop:window} is worth stating because it says which
experimental quantity matters and which does not. Total OSK expression level,
raised or lowered uniformly, does not move the window; it moves the time
scale. What moves the window is the \emph{ratio} of propensities at the
contested sites. An intervention that increases OSK abundance while leaving
that ratio alone buys speed and no safety, which is a prediction that
contradicts the intuition that lower dose is safer.

\begin{remark}[Where the intuition is right anyway]\label{rem:window2}
Lower total exposure does reduce risk, but by a different route: it reduces
the number of firings, hence the number of opportunities for a low-probability
$\Ded$ event, and it shortens the interval during which identity is
destabilised. That is a statement about the \emph{number of trials}, not about
the per-trial odds, and the model should keep the two separate because they
respond to different design levers.
\end{remark}

\subsection{What the cell's threshold is}

Both thresholds require a definition in the term, and neither is obvious. The
rejuvenation threshold is a predicate on the methylation configuration of the
regeneration-competence loci --- a formula in the crisp sort, of the kind
Layer~1's keys already use. The dedifferentiation threshold is harder, because
loss of identity is not a single locus event; it is plausibly a fixed-point
property of the pluripotency network's mutual activation, which lands it in
$\mathcal{F}_3$ and hence, by Principle~4.6 of \cite{Interventions}, in the
model checker rather than in the clause.

\begin{obligation}[Thresholds as formulae]\label{ob:thresh}
Give both thresholds as formulae in the generated logic over the chromatin
term, with the dedifferentiation threshold stated in $\mathcal{F}_3$ and
lowered to $\mathcal{F}_1$ on the finite-state fragment by unfolding, per
Definition~4.4 of \cite{Interventions}. Until this is done, the simulator has
no stopping condition and the window has no definition.
\end{obligation}

%%%%%%%%%%%%%%%%%%%%%%%%%%%%%%%%%%%%%%%%%%%%%%%%%%%%%%%%%%%%%%%%%%%%%%%%%%%%%%%
\section{Fidelity: a proofreading bound and a cheap test}\label{sec:fidelity}

\subsection{Pulsing is proofreading in time}

OSK must bind the right loci. The discrimination is between binding sites that
differ by a small free-energy margin, which is the classical setting of
Hopfield--Ninio kinetic proofreading. The companion note establishes the
relevant structure: a catalytic step is a context rule, restoring its
left-hand side and contributing a multiplicative factor at no primitive cost,
while a proofreading step is a base rule, consuming its redex and clearing the
funding judgement \cite{Proofreading}. Selectivity beyond the binding energies
requires the second kind.

A pulsed induction regime has exactly that shape. Each withdrawal phase
releases complexes that have not committed, which is a discard step; each
subsequent induction phase re-draws. The ladder's rungs are pulses and its
currency is expression.

\begin{observation}[The exchange rate]\label{obs:rate}
By the exchange rate of \cite{Proofreading}, error falls as a power of the
resource burned, with an exponent determined by single-stage kinetics and
independent of the depth of the ladder. Writing $f$ for the single-stage
discrimination factor at the contested loci --- the ratio of off-rates for
incorrect versus correct occupancy --- the achievable error after $n$ discard
stages scales as $f^{-n}$ in the ideal limit, with the realised exponent
degraded by re-binding during the off phase.
\end{observation}

\begin{prediction}[Pulsed versus sustained]\label{pred:pulse}
Pulsed induction improves selectivity over sustained induction if and only if
$f > 1$ at the contested loci, and the magnitude of the improvement is
governed by $f$ and by the ratio of the off-phase duration to the re-binding
time, not by the number of pulses beyond the point at which re-binding
dominates. If $f \approx 1$, schedule optimisation is not the design problem.
\end{prediction}

Prediction~\ref{pred:pulse} is ordinal and falsifiable, and it does not
require any of the unknown rate constants of \S\ref{sec:gaps}: it requires
$f$, which is a two-parameter binding measurement.

\subsection{The limitative result, and what it implies for the programme}

The sharper content is negative.

\begin{principle}[Energy buys confidence and never scope]\label{prin:scope}
A proofreading step is a move in the fibre \cite{Proofreading}. It can
increase the reliability with which a discrimination already available is
made; it cannot create a discrimination that the affinity landscape does not
support. Consequently no induction schedule, however elaborate, manufactures
selectivity between regeneration-competence loci and pluripotency loci if OSK
does not distinguish them.
\end{principle}

The programmatic consequence is a fork, and it is worth facing before any
animal work.

\begin{enumerate}[leftmargin=1.6em,itemsep=3pt]
\item If $f$ is comfortably greater than one at the contested loci, schedule
design is the right lever, Prediction~\ref{pred:pulse} tells you which
direction to push, and the model's job is to find the schedule that maximises
the window subject to the trial-count consideration of
Remark~\ref{rem:window2}.
\item If $f$ is near one, no schedule helps, and the design problem is a
different one: a different factor set, or a targeting layer that supplies the
discrimination exogenously --- locus-directed recruitment rather than diffuse
overexpression. The model would then be a model of the wrong intervention, and
the right response is to change the intervention rather than to tune it.
\end{enumerate}

\begin{principle}[The go/no-go precedes the model]\label{prin:gonogo}
Measuring $f$ on hPSC-derived hindbrain motor neurons \cite{TwoOrgans} is a
binding assay on the correct substrate. It is cheaper than a cohort, it is
upstream of every parameter the model lacks, and by
Principle~\ref{prin:scope} it determines which of the two programmes above is
the one being pursued. It should therefore be the first bench work, before
model fitting and before animal work.
\end{principle}

This is the note's most actionable finding, and it is a consequence of the
formalism rather than of the biology: the limitative theorem is what tells you
that this particular measurement screens off the rest.

%%%%%%%%%%%%%%%%%%%%%%%%%%%%%%%%%%%%%%%%%%%%%%%%%%%%%%%%%%%%%%%%%%%%%%%%%%%%%%%
\section{Counterfactual cohort design}\label{sec:cohort}

Rung three of the ladder has a concrete payoff here, and it is statistical
rather than philosophical.

A run of the model is a causal history together with the candidate that fired
at each site; its exogenous record is the vector of resolver draws, and a run
is replayed under $\dok(C)$ by holding that record and re-resolving only where
the restriction bites \cite[Def.~3.12, Thm.~3.13]{Interventions}. For a
simulated cell this answers: \emph{would this cell have degenerated without
the intervention?} --- with the same noise realisation on both arms.

\begin{observation}[Matched pairs for free]\label{obs:pairs}
Replay gives a per-cell matched-pair contrast rather than a difference of
cohort means. The variance of the estimated treatment effect is accordingly
the variance of the within-pair difference, which removes the between-cell
variation that dominates a two-cohort comparison of heterogeneous neurons.
\end{observation}

Two cautions. First, the clean statement of Theorem~3.13 of
\cite{Interventions} holds on the staged fragment $\mathcal{T}_1$, where every
gadget's contention set is itself and the recorded class is always still a
candidate; a well-mixed cell is not staged, so off that fragment the ``and
otherwise sample'' branch is taken and the matching degrades. Quantifying how
much it degrades, for this model, is Obligation~\ref{ob:replay}. Second, the
variance reduction is a property of the \emph{simulation}, not of a real
cohort: no animal study has access to the exogenous record. The payoff is in
designing the study --- power calculations, which readouts at which timepoints
--- not in analysing it.

\begin{obligation}[Replay fidelity off the staged fragment]\label{ob:replay}
Measure, in simulation, the fraction of replayed firings at which the recorded
candidate is no longer available, as a function of the intervention's position
in the cascade. This bounds the usefulness of counterfactual contrasts for
this model and is cheap to obtain once Layer~0 runs.
\end{obligation}

%%%%%%%%%%%%%%%%%%%%%%%%%%%%%%%%%%%%%%%%%%%%%%%%%%%%%%%%%%%%%%%%%%%%%%%%%%%%%%%
\section{Four places the apparatus is not ready}\label{sec:gaps}

Each of these is a real gap. Each has a workaround with a stated cost. None is
fatal to the programme, and two of them bound what the model may claim.

\subsection{Cycles at variable level}

The pluripotency network is mutually activating: Oct4, Sox2 and Nanog
reinforce one another, and the autoregulation is the mechanism by which
dedifferentiation, once begun, completes. So the variable-level graph has
cycles, while Pearl's apparatus and the surgery theorem of
\cite{Interventions} are stated where the graph is a DAG.

The distinction the companion note draws is the one to use: the
\emph{event-level} graph is always acyclic, because an event structure's
configurations are partial orders; the \emph{variable-level} graph, which
identifies events of a kind, may cycle when the term recurses. The workaround
is to work at event level with unrolling depth equal to the inference budget
$\beta$, treating the variable-level cycle as its unrolling.

The cost is that every interventional claim is indexed by $\beta$, and the
$\beta$-dependence must be reported rather than suppressed. A claim that
stabilises in $\beta$ is trustworthy; one that does not is an artefact of the
truncation. Obligation~5.6 of \cite{Interventions} --- cyclic SCMs meeting the
graded fixed point \cite{BongersEtAl} --- is the principled route and is open.

\subsection{Well-mixed contention sets}

A cell is a beaker. There are many copies of each species, so contention sets
are large, and the normaliser is a sum over maximal matchings, which is
permanent-shaped and $\#\mathsf{P}$-hard in general
\cite[\S4.1]{Interventions}. Discipline $\mathcal{T}_1$ buys singleton
contention sets, but a cell does not satisfy it.

The architecture of Table~\ref{tab:layers} is the workaround. The Gillespie
propensity reading handles multiplicity directly, as the $k=1$ truncation ---
the ideal gas of \cite{ProcessMolecules} --- where propensity is a count times
a rate and no permanent arises. The $\dok$ membrane with $\mathcal{F}_0$ or
$\mathcal{F}_1$ constraints is confined to Layer~2's designated intervention
sites, which are few and whose contention sets are small by construction.

\begin{principle}[Do not put the cell under a $\dok$]\label{prin:nodok}
The membrane is scoped to the sites the intervention acts on. Bulk chemistry
runs beneath it. A $\dok$ whose scope is the whole cell forces recomputation
of large contention-set normalisers and is the proximate cause of
intractability in models of this shape.
\end{principle}

Spatial structure --- nuclear compartmentalisation, condensate formation ---
would require the collision kernel of \cite{ProcessMolecules} and its higher
virial terms. That machinery should be deferred: it is not needed for the
questions of \S\ref{sec:race} and \S\ref{sec:fidelity}, and importing it early
would cost the identifiability of \S\ref{sec:gaps}.3 for no gain.

\subsection{Identifiability, which is the real risk}

Most of the rate constants are unknown, and this is where the project fails if
it fails. The gradient machinery is available: over $\Rnn$ a modality-free
clause is a posynomial in its parameters, hence smooth and log-convex in
log-coordinates, the likelihood factors over firing sites, and the parameter
gradient decomposes into local scores computable at each site from candidate
values the site already had \cite{GradedPLN}. Fitting is therefore
mechanically straightforward.

Fitting is not identification.

\begin{principle}[Identifiability analysis precedes estimation]
\label{prin:ident}
A model with many unidentifiable parameters that reproduces the available
data is a curve, not a model: it will fit, it will interpolate, and its
interventional predictions --- the only predictions wanted here --- will be
determined by the prior rather than the data. Structural identifiability
analysis therefore precedes parameter estimation, and its output is a ranked
list of candidate measurements by information gain about the quantities
\S\ref{sec:race} and \S\ref{sec:fidelity} actually use.
\end{principle}

The sharp form of the question is Obligation~5.4 of \cite{Interventions}:
Pearl's identifiability problem rendered as a question about which
interventional quantities are recoverable from the observational data one can
actually obtain. Here the observables are realistic and limited --- single-cell
transcriptomes at a few timepoints, methylation arrays, immunostaining,
occupancy assays --- and the hidden variables are the states no assay reports.
Transferring the ID algorithm \cite{ShpitserPearl} to the staged fragment is
the stated route and is undischarged.

\begin{obligation}[Identifiability before fitting]\label{ob:ident}
Perform structural identifiability analysis on the Layer~0 and Layer~1 model
before estimating any parameter, and report which of
Proposition~\ref{prop:window}'s ratio and Observation~\ref{obs:rate}'s
exponent are identifiable from which observable sets. This is the deliverable
on which the model's usefulness depends.
\end{obligation}

\subsection{Starvation discipline}

A restricted receipt with no satisfying candidate waits forever: the history
deadlocks by starvation, which is how the embedding renders the failure of
Pearl's positivity condition \cite[Rem.~3.8]{Interventions}. In a biological
model a starvation may be a real dead end --- a cell state from which the
cascade cannot proceed --- or an artefact of an over-restricted $\dok$. These
must not be confused, and in simulation they look identical.

Under discipline $\mathcal{T}_3$, where every channel a constrained variable
binds on carries a type with finitely many closed values, vacuity and fatality
are decidable by inspecting which values the clause weights positively
\cite[Def.~4.5]{Interventions}. The discipline to adopt is accordingly:

\begin{principle}[Every $\dok$ ships with its vacuity check]\label{prin:vac}
Each intervention term in the model carries a decidable vacuity-and-fatality
check, discharged statically under $\mathcal{T}_3$. Any starvation observed in
simulation is triaged against that check before being interpreted
biologically.
\end{principle}

%%%%%%%%%%%%%%%%%%%%%%%%%%%%%%%%%%%%%%%%%%%%%%%%%%%%%%%%%%%%%%%%%%%%%%%%%%%%%%%
\section{What the model can and cannot license}\label{sec:license}

This section is the one to agree before bench work, because it is the
boundary that makes the rest usable.

\subsection{Licensed}

\begin{itemize}[leftmargin=1.4em,itemsep=3pt]
\item \textbf{Ranking schedules.} Ordinal comparison of induction regimes by
the window of Proposition~\ref{prop:window}, under stated kinetics, with the
sensitivity of the ranking reported alongside it.
\item \textbf{Prioritising measurements.} Identification of which parameters
the ranking depends on, hence which assays buy the most information per unit
cost (Obligation~\ref{ob:ident}).
\item \textbf{Falsifiable ordinal predictions.} Prediction~\ref{pred:pulse} is
testable without knowing the unknown rates.
\item \textbf{Discarding variants.} A protocol variant that the model shows to
be dominated under every parameterisation consistent with the data need not be
tested. Negative screening is where a model of this kind earns its keep.
\item \textbf{Detecting structural problems early.} Non-identifiability and
transfer obstructions (Proposition~\ref{prop:wf}) are found statically,
before money is spent.
\item \textbf{Study design.} Counterfactual contrasts for power calculation
and readout timing (\S\ref{sec:cohort}).
\end{itemize}

\subsection{Not licensed}

\begin{itemize}[leftmargin=1.4em,itemsep=3pt]
\item \textbf{Safety.} The model cannot establish it. The events that matter
--- transformation, loss of identity in a critical nucleus --- are rare, their
governing parameters are the least known in the model, and rare-event
probabilities are precisely where a model with unidentified parameters is
least trustworthy. A simulated window is not a safety margin.
\item \textbf{Magnitudes of oncogenic risk.} Withholding c-Myc reduces a risk
that the model cannot quantify, because quantifying it requires the tail of a
distribution the model does not constrain.
\item \textbf{Substituting for toxicology or animal work.} The model reorders
and prunes that work; it does not replace any of it.
\item \textbf{Any dose, route, or schedule for administration to a person.}
Out of scope by \S\ref{sec:scope}.3 and not recoverable from anything here.
\end{itemize}

\begin{principle}[What the model is for]\label{prin:for}
The model's job is to say which experiments to run and which protocol variants
to discard. It is not to say which protocol to use.
\end{principle}

%%%%%%%%%%%%%%%%%%%%%%%%%%%%%%%%%%%%%%%%%%%%%%%%%%%%%%%%%%%%%%%%%%%%%%%%%%%%%%%
\section{A staged programme}\label{sec:programme}

The ordering below is chosen so that each stage produces something checkable
and so that the stages most likely to invalidate the approach come early.

\begin{enumerate}[leftmargin=1.6em,itemsep=4pt]
\item[\textbf{S0}] \textbf{Substrate and populations.} Fix the bulbar motor
pool and the off-target population of Principle~\ref{prin:offtarget};
enumerate species and loci; state both thresholds
(Obligation~\ref{ob:thresh}).
\item[\textbf{S1}] \textbf{Layer 0, as a SPiM model.} Static maps,
channel-identity keys. Literature rates where they exist, explicit
\textsf{unknown} flags where they do not. Readable and runnable outside this
programme by Remark~\ref{rem:spim}.
\item[\textbf{S2}] \textbf{Layer 1, the dynamic-map lift.} The first
publishable unit; see below.
\item[\textbf{S3}] \textbf{Layer 2, the intervention.} Schedule reified per
Proposition~\ref{prop:reify}; tropism as $\mathcal{T}_4$; schedules as terms.
\item[\textbf{S4}] \textbf{Layer 3 and the $f$ measurement.} The go/no-go of
Principle~\ref{prin:gonogo}. \emph{This is the first bench work.}
\item[\textbf{S5}] \textbf{Identifiability and sensitivity.}
Obligation~\ref{ob:ident}. The deliverable for experimental collaborators.
\item[\textbf{S6}] \textbf{Cohort design.} \S\ref{sec:cohort}, conditional on
Obligation~\ref{ob:replay}.
\end{enumerate}

Two remarks on the ordering. The $f$ measurement (S4) is placed before
identifiability and fitting because Principle~\ref{prin:scope} says it screens
off which programme is being pursued; spending on model calibration before
knowing the answer risks calibrating a model of the wrong intervention. And
S1 is retained as a deliverable in its own right rather than folded into S2,
because an externally checkable artefact is worth more than the week it costs.

\subsection{The first publishable unit}

S2 is the natural first paper: \emph{dynamic weight maps for epigenetic state,
with a reprogramming cascade as the worked example}. It is a formal result
about expressiveness --- Observation~\ref{obs:blowup}, with the SPiM
restriction theorem as its baseline --- and so does not depend on the unknown
rate constants. It is also the piece that makes the transfer obstruction of
\S\ref{sec:transfer} statable, since that argument is about key sets.

The transferability argument itself is a plausible second, and a shorter one:
\emph{why a model calibrated on anterior-lineage CNS neurons does not transfer
to posterior-lineage motor neurons}. It would be read by a different audience
than this programme's usual one, it makes a prediction the hPSC substrate of
\cite{TwoOrgans} can test directly, and
Proposition~\ref{prop:wf} gives it a sharper form than the biological
literature currently has.

%%%%%%%%%%%%%%%%%%%%%%%%%%%%%%%%%%%%%%%%%%%%%%%%%%%%%%%%%%%%%%%%%%%%%%%%%%%%%%%
\section{Literature points to settle first}\label{sec:settle}

Four factual matters bear on the model's construction and should be confirmed
against the primary sources rather than carried from summaries.

\begin{enumerate}[leftmargin=1.6em,itemsep=3pt]
\item \textbf{Switch polarity.} The system of \cite{LuEtAl} uses tTA under a
constitutive promoter with TRE-OSK, and the primary material describes
induction in the \emph{absence} of doxycycline, which is Tet-Off. Secondary
summaries state the opposite. The polarity determines the sign of the
schedule term in Layer~2 and must be taken from the primary source.
\item \textbf{TET1/TET2 dependence across lineages.} The requirement for the
demethylases is established on anterior substrate. Whether it holds, and with
what stoichiometry, in posterior-lineage motor neurons is a parameter with a
lineage index, and by Observation~\ref{obs:keys} possibly a key with one.
\item \textbf{Which motor pools are in scope.} Spinal motor neurons derive
from neuromesodermal progenitors rather than from hindbrain neural ectoderm,
so the hindbrain constraint restricts the target to the cranial motor nuclei.
If the intent is to include spinal pools, that is a third lineage and a third
partition, not an extension of the second.
\item \textbf{Current clinical status of the anterior programme.} An
OSK-based therapy derived from this work has entered human trials for
glaucoma. What has been published about its dosing and safety profile is
directly relevant to Layer~2 and should be tracked, with the caveat that it is
an anterior-organ, acute-accessible-tissue indication.
\end{enumerate}

%%%%%%%%%%%%%%%%%%%%%%%%%%%%%%%%%%%%%%%%%%%%%%%%%%%%%%%%%%%%%%%%%%%%%%%%%%%%%%%
\section{Conclusion}\label{sec:conc}

The apparatus fits this problem better than it had any particular right to,
and the fit is not coincidental: a therapeutic window is a statement about
which of two competing futures a cell enters, and that is the structure the
graded where-clause was built to make writable. Rejuvenation and
dedifferentiation contend for one OSK pool, so the window is a property of the
contention set's normalisation rather than of either propensity alone, and
raising expression uniformly buys time rather than safety. The induction
schedule must be reified as data before any intervention operator can reach
it. Epigenetic state dependence needs dynamic weight maps, which is where the
construction leaves SPiM behind, and SPiM remains available underneath as the
layer collaborators can check.

Two results point outward rather than inward. The published calibration does
not transfer across the organ boundary, and the partition discipline makes
that a well-formedness question rather than a matter of opinion --- with the
persistence of progenitor divergence into maturity as the empirical assumption
the argument rests on. And the proofreading ladder, read for its limitative
content, says that no schedule creates selectivity the affinity landscape does
not supply, which converts a theorem into the cheapest consequential
experiment available: measure the single-stage discrimination factor on
hindbrain motor neurons, and let it decide which of two different programmes
this is.

The model will not establish that anything is safe. What it can do is spend
the next year's experiments in a better order than intuition would, and say
plainly which of its own conclusions are load-bearing and which are furniture.
That is the whole of what is on offer, and it is worth having.

%%%%%%%%%%%%%%%%%%%%%%%%%%%%%%%%%%%%%%%%%%%%%%%%%%%%%%%%%%%%%%%%%%%%%%%%%%%%%%%
\begin{thebibliography}{99}\raggedright

\bibitem{LuEtAl}
Y.~Lu, B.~Brommer, X.~Tian, A.~Krishnan, M.~Meer, C.~Wang, D.~L.~Vera,
Q.~Zeng, D.~Yu, M.~S.~Bonkowski, J.-H.~Yang, S.~Zhou, E.~M.~Hoffmann,
M.~M.~Karg, M.~B.~Schultz, A.~E.~Kane, N.~Davidsohn, E.~Korobkina,
K.~Chwalek, L.~A.~Rajman, G.~M.~Church, K.~Hochedlinger, V.~N.~Gladyshev,
S.~Horvath, M.~E.~Levine, M.~S.~Gregory-Ksander, B.~R.~Ksander, Z.~He, and
D.~A.~Sinclair.
\newblock Reprogramming to recover youthful epigenetic information and restore
vision.
\newblock \emph{Nature} 588:124--129, 2020.
\newblock \url{https://doi.org/10.1038/s41586-020-2975-4}

\bibitem{TwoOrgans}
R.~Jokhai, L.~Dundes, et al.\ (K.~M.~Loh, senior author).
\newblock Two parallel lineage-committed progenitors contribute to the
developing brain.
\newblock \emph{Nature Neuroscience}, published 18 September 2026.
\newblock Preprint: \emph{bioRxiv} 2025.07.02.662771.
\newblock \url{https://www.ncbi.nlm.nih.gov/pmc/articles/PMC12236623/}

\bibitem{OcampoEtAl}
A.~Ocampo, P.~Reddy, P.~Martinez-Redondo, A.~Platero-Luengo, F.~Hatanaka,
T.~Hishida, et al.
\newblock In vivo amelioration of age-associated hallmarks by partial
reprogramming.
\newblock \emph{Cell} 167(7):1719--1733, 2016.

\bibitem{PhillipsCardelli}
A.~Phillips and L.~Cardelli.
\newblock Efficient, correct simulation of biological processes in the
stochastic pi-calculus.
\newblock In \emph{Computational Methods in Systems Biology}, LNCS 4695,
pp.~184--199, 2007.

\bibitem{PhillipsCardelliMachine}
A.~Phillips and L.~Cardelli.
\newblock A correct abstract machine for the stochastic pi-calculus.
\newblock In \emph{Concurrent Models in Molecular Biology (BioCONCUR)}, 2004.

\bibitem{Gillespie}
D.~T.~Gillespie.
\newblock Exact stochastic simulation of coupled chemical reactions.
\newblock \emph{Journal of Physical Chemistry} 81(25):2340--2361, 1977.

\bibitem{Hopfield}
J.~J.~Hopfield.
\newblock Kinetic proofreading: a new mechanism for reducing errors in
biosynthetic processes requiring high specificity.
\newblock \emph{PNAS} 71(10):4135--4139, 1974.

\bibitem{Ninio}
J.~Ninio.
\newblock Kinetic amplification of enzyme discrimination.
\newblock \emph{Biochimie} 57(5):587--595, 1975.

\bibitem{Pearl09}
J.~Pearl.
\newblock \emph{Causality: Models, Reasoning, and Inference}.
\newblock Cambridge University Press, second edition, 2009.

\bibitem{ShpitserPearl}
I.~Shpitser and J.~Pearl.
\newblock Identification of joint interventional distributions in recursive
semi-Markovian causal models.
\newblock In \emph{AAAI 2006}, pp.~1219--1226.

\bibitem{BongersEtAl}
S.~Bongers, P.~Forr\'e, J.~Peters, and J.~M.~Mooij.
\newblock Foundations of structural causal models with cycles and latent
variables.
\newblock \emph{Annals of Statistics} 49(5):2885--2915, 2021.

\bibitem{MeredithRadestock}
L.~G.~Meredith and M.~Radestock.
\newblock A reflective higher-order calculus.
\newblock \emph{Electronic Notes in Theoretical Computer Science}
141(5):49--67, 2005.

\bibitem{Interventions}
L.~G.~Meredith.
\newblock Interventions at the cut: Pearl's \textsf{do} as a scoped guard in
the rho calculus with graded where-clauses.
\newblock F1R3FLY.io working note,
\texttt{fuzzyware/interventions-at-the-cut}, November 2026.

\bibitem{GradedV4}
L.~G.~Meredith.
\newblock Graded where-clauses: non-crisp truth values as the syntax of
non-determinism resolution in the rho calculus, version 4.
\newblock F1R3FLY.io working note, \texttt{fuzzyware/graded-where-v4}, August
2026.

\bibitem{GradedPLN}
L.~G.~Meredith.
\newblock Graded where-clauses over the reals: one-dimensional grading as the
mortal scientist's gradient, and probabilistic logic as the metric that
gradient requires.
\newblock F1R3FLY.io working note, \texttt{fuzzyware/graded-where-pln}, August
2026.

\bibitem{WeightedGSLT}
L.~G.~Meredith.
\newblock Weighted GSLTs: stochastic and quantum simulation for decorated
continued interactive GSLTs.
\newblock F1R3FLY.io working note, \texttt{rho-life/weighted-gslt}, 2026.

\bibitem{Proofreading}
L.~G.~Meredith.
\newblock Kinetic proofreading in the mortal scientist setting: energy into
accuracy, and the price of commitment.
\newblock F1R3FLY.io working note,
\texttt{kinetic-proofreading/proofreading-note}, August 2026.

\bibitem{ProcessMolecules}
L.~G.~Meredith.
\newblock Process terms as molecules: locations, collisions, and the road to
viscosity.
\newblock F1R3FLY.io working note, \texttt{fuzzyware/process-molecules}, 2026.

\bibitem{Inhabiting}
L.~G.~Meredith.
\newblock Inhabiting choice: what a confluent witness cannot be.
\newblock F1R3FLY.io working note,
\texttt{proof-dynamics-and-reasoning/inhabiting-choice}, November 2026.

\bibitem{ProofSearch}
L.~G.~Meredith.
\newblock Proof search as concurrent reaction.
\newblock F1R3FLY.io working note,
\texttt{proof-dynamics-and-reasoning/proof-search}, 2026.

\bibitem{CostMonad}
L.~G.~Meredith.
\newblock Cost accounting as an endofunctor on continued interactive GSLTs.
\newblock F1R3FLY.io working note,
\texttt{cost-accounting-as-monad/continued-gslt-cost-v2}, 2026.

\end{thebibliography}

\end{document}
